\section{Kagero UI引擎}
\label{sec:kagero}

\subsection{背景}

游戏UI系统需要处理复杂的用户交互、动态布局和数据绑定。现代游戏UI的特点包括：

\begin{enumerate}
    \item \textbf{状态驱动}：UI需要响应游戏状态的变化（如玩家血量、物品栏变化）
    \item \textbf{布局复杂}：需要支持多种布局模式（流式、网格、锚点等）
    \item \textbf{事件多样}：鼠标点击、键盘输入、焦点变化等
    \item \textbf{性能敏感}：UI更新不应阻塞游戏逻辑
\end{enumerate}

传统立即模式GUI在复杂场景下性能不佳，而保留模式GUI框架往往过于重量级，不适合游戏场景。

\subsection{原版局限性}

Java版Minecraft的GUI系统存在以下问题：

\begin{enumerate}
    \item \textbf{立即模式渲染}：每帧重建所有绘制命令，无法有效复用
    \item \textbf{无数据绑定}：状态变化需要手动更新UI，容易遗漏或不一致
    \item \textbf{布局硬编码}：位置和大小通常写死在代码中，难以适应不同分辨率
    \item \textbf{无模板系统}：相似的UI组件重复实现，代码冗余
    \item \textbf{状态同步困难}：游戏状态和UI状态分散在多处，维护成本高
\end{enumerate}

\subsection{本项目实现}

Kagero（陽炎）是自研的响应式UI引擎，采用模板编译和状态绑定设计，命名为Kagero寓意UI如热浪般流畅自然。

\subsubsection{响应式状态管理}

核心是\texttt{Reactive<T>}模板类，实现观察者模式的自动状态追踪：

\textbf{基本原理}：每个响应式值维护一个观察者列表。当值变化时（通过\texttt{set}方法），自动通知所有观察者。观察者可以是一个回调函数，用于更新UI。

\textbf{值比较优化}：在通知前比较新旧值，只有真正变化才触发通知，避免无效更新。

\textbf{计算属性}：\texttt{Computed<T>}允许定义依赖其他响应式值的计算属性，当依赖变化时自动重新计算。

\textbf{批量更新}：\texttt{StateStore}支持批量更新，在批量更新期间收集所有变化的键，更新结束后统一通知订阅者。这避免了多次状态变化导致的多次UI重绘。

\subsubsection{模板编译系统}

Kagero使用XML风格的模板语言描述UI结构：

\textbf{编译流程}：模板源码经过词法分析（Lexer）生成Token流，再经过语法分析（Parser）生成抽象语法树（AST），最后由模板编译器（TemplateCompiler）生成可执行的CompiledTemplate。

\textbf{AST节点类型}：Document（文档根节点）、Screen（屏幕节点）、Widget（通用组件）、Button、Text、TextField等。每个节点包含静态属性、绑定属性和事件属性。

\textbf{绑定上下文}：\texttt{BindingContext}负责将模板中的绑定表达式解析为实际值。支持：
\begin{itemize}
    \item 暴露C++变量供模板绑定（只读和可写）
    \item 暴露Reactive状态
    \item 暴露回调函数
\end{itemize}

\textbf{循环渲染}：使用\texttt{for:item="item in items"}语法循环渲染列表。循环变量在子作用域中有效。

\textbf{条件渲染}：使用\texttt{if="condition"}语法条件渲染组件。

\subsubsection{布局引擎}

支持多种布局算法，通过适配器模式将Widget适配为布局节点：

\textbf{测量规格}：包含大小和模式两个维度。模式包括：
\begin{itemize}
    \item \texttt{Unspecified}：无约束，子元素可以使用任意大小
    \item \texttt{Exactly}：精确值，子元素必须使用指定大小
    \item \texttt{AtMost}：最大值，子元素可以不超过指定大小
\end{itemize}

测量规格的解析：
\begin{equation}
size_{measured} = \begin{cases}
size_{spec} & \text{if mode = Exactly} \\
\min(size_{desired}, size_{spec}) & \text{if mode = AtMost} \\
size_{desired} & \text{if mode = Unspecified}
\end{cases}
\end{equation}

\textbf{Flex布局}：支持主轴方向（行/列）、主轴对齐（start/center/end/space-between/space-around）、交叉轴对齐（start/center/end/stretch）、换行模式等配置。

\textbf{布局过程}：两遍测量，第一遍测量所有子元素，第二遍根据测量结果分配位置。

\subsubsection{事件系统}

事件总线采用发布-订阅模式，支持类型安全的事件分发：

\textbf{事件类型}：定义了50+种事件类型，包括鼠标事件（Click、Release、Drag、Scroll、Move、Enter、Leave）、键盘事件（KeyPress、KeyRelease、KeyRepeat、CharInput）、焦点事件（FocusGained、FocusLost）、值变化事件（ValueChange、TextChange）、Widget事件（Init、Resize、Show、Hide）等。

\textbf{优先级}：订阅者可以指定优先级，高优先级的订阅者优先处理事件。相同优先级按注册顺序处理。

\textbf{事件过滤器}：支持在事件分发前进行过滤，过滤器可以阻止事件继续传播。

\textbf{RAII订阅}：\texttt{EventSubscription}类在析构时自动取消订阅，避免内存泄漏。

\subsubsection{渲染抽象层}

\texttt{ICanvas}接口定义了平台无关的绘图操作：

\textbf{基本图形}：矩形、圆角矩形、圆形、椭圆、路径、直线。

\textbf{渐变}：线性渐变，支持垂直和水平方向。

\textbf{图像}：图像绘制、九宫格拉伸（Nine-Patch）。

\textbf{文本}：文本绘制、文本测量。

\textbf{裁剪和变换}：矩形裁剪、圆角矩形裁剪、路径裁剪、平移、缩放、旋转、矩阵变换。

\textbf{状态栈}：save/restore机制，保存和恢复当前状态（裁剪区域、变换矩阵）。

\texttt{PaintContext}在ICanvas之上提供便捷方法，如居中文本绘制、边框绘制、填充矩形等，并缓存常用的Paint对象。

\subsection{解决与成果}

\begin{table}[h]
\centering
\caption{UI系统对比}
\label{tab:ui_comparison}
\begin{tabularx}{\linewidth}{Xcc}
\toprule
\textbf{特性} & \textbf{Java版} & \textbf{Kagero} \\
\midrule
渲染模式 & 立即模式 & 保留模式 \\
数据绑定 & 无 & Reactive<T> \\
布局系统 & 硬编码 & Flex/Grid/Anchor \\
模板系统 & 无 & XML模板编译 \\
事件处理 & 回调注册 & EventBus \\
性能优化 & 无 & 增量布局、批量更新 \\
\bottomrule
\end{tabularx}
\end{table}

通过模板编译和状态绑定，UI代码量减少约60\%，状态同步错误减少约90\%。响应式设计使得UI自动适应游戏状态变化，开发者无需手动管理状态同步。布局引擎支持动态调整大小，适应不同分辨率和窗口大小。
